\section{Preparation and selection proofs}
\label{sec:supplementary-selection-proofs}

\subsection{Existence of admissible preparations}
\label{subsec:supplementary-preparation-existence}

We first record the implementation hidden behind the concise selection
convention in the main article.  For the graph rule, all state derivatives
and mixed parameter derivatives
\begin{equation}
 D_u^a\partial_\rho^b\partial_\nu^i\partial_\zeta^j,
 \qquad a+b+i+j\le12,\qquad b\le3,\qquad i\le2,\qquad j\le2,
                                                               \label{eq:preparation-derivative-list}
\end{equation}
are controlled in the product norms.  Here \(\rho\) is the dummy chart
amplitude used in the graph proof.  The rule uses one scalar cutoff and one
stable-saturation profile for all \(N\), applies the saturation
componentwise followed by \(P_{\perp,N}\), and is independent of
\((\nu,\zeta)\).

Choose \(\omega\in C_c^\infty(\R;[0,1])\), equal to one on \([-1,1]\) and
zero off \((-2,2)\), and constants
\(B_*>2\Theta_*+B+3\) and \(d_*>b_0\) such that, for every sufficiently
small target, the depth-two agreement hull lies in
\(\{|\chi|<R_\delta+B_*,\ |d|<d_*\}\).  The frozen singular field is
\begin{equation}
 \Xi_\delta(\chi,d)
 \binom{1-2\alpha d}{\chi d},\qquad
 \Xi_\delta(\chi,d)=
 \omega\!\left(\frac{\chi}{R_\delta+B_*}\right)
 \omega\!\left(\frac d{d_*}\right).                 \label{eq:fixed-singular-preparation}
\end{equation}
The graph transform uses the \((\chi,d)\)-norm.  Its frozen preparation is
tame through grade eleven in the sense of
\cref{def:invdet-tame-preparation}; in particular,
\begin{equation}
 \operatorname{Lip}_{\chi,d}(q_{0,\delta})
 \le C_{\rm prep}(1+R_\delta).                       \label{eq:fixed-singular-lipschitz}
\end{equation}

The planar completion uses precisely the continuous mixed jets
\begin{equation}
 D_u^a\partial_\nu^i\partial_\zeta^j,
 \qquad 0\le a\le3,\quad 0\le i,j\le2,\quad i+j\le2. \label{eq:planar-required-jets}
\end{equation}
It is globally bounded and Lipschitz, agrees with the reduced field on the
retained tube, equals \(q_0\) on the two endpoint balls, and obeys
\cref{eq:planar-preparation-smallness}.  Both extension rules remain fixed
when parameter derivatives are taken.  The smallness condition excludes
order-one forcing in an artificial transition strip while retaining the
physical \(O(\delta)\) term and its \(\nu\)-derivative, together with the
\(O(\delta^2)\) structural and second-parameter sources used by the trace
solve.

We call a graph rule \(\mathfrak P_{\rm gr}\) \emph{graph-admissible} if it
satisfies clauses \textup{(S1)}--\textup{(S2)} of
\cref{def:admissible-preparation}, including
\cref{eq:fixed-singular-preparation,eq:fixed-singular-lipschitz} and the
grade-eleven tame-flow requirement, without presupposing a planar completion.

\begin{lemma}[Graph-only retained-tube expansion]
\label[lemma]{lem:graph-only-retained-expansion}
Fix the common model bounds and a graph-admissible rule
\(\mathfrak P_{\rm gr}\).  After decreasing a common target threshold, the
special-flow transform has a unique invariant history graph, with every
mixed jet indexed by \(\mathcal I_{10}\) in
\cref{eq:invdet-index-set}; in particular, it has the full rectangular jet
family in \cref{eq:invdet-required-rectangle}.  On the retained set its
reduced field has the exact dummy-amplitude expansion
\[
 Q_{N,\delta,\nu,\zeta}
 =q_0+\delta q_{1,N}+\delta^2q_{2,N}
       +\delta^3\mathcal R_{3,N}.
\]
There are \(C,c>0\) and \(M\in\mathbb N\), independent of \(N\) and
\(\delta\), such that, for \(0\le a\le3\), \(0\le i,j\le2\), and
\(i+j\le2\),
\begin{equation}
 \max\left\{
 |D_u^a\partial_\nu^i\partial_\zeta^j q_{1,N}|,
 |D_u^a\partial_\nu^i\partial_\zeta^j q_{2,N}|,
 |D_u^a\partial_\nu^i\partial_\zeta^j\mathcal R_{3,N}|
 \right\}(\Gamma(\chi,d))
 \le C e^{c|\chi|}\langle\chi\rangle^M.             \label{eq:graph-only-retained-envelope}
\end{equation}
throughout
\[
 |\chi|\le S_\delta+1,\qquad |d|\le\delta^\kappa,\qquad
 \nu\in I_\nu,\qquad |\zeta|\le\zeta_*.
\]
Moreover
\(\partial_\zeta q_{1,N}=\partial_{\nu\nu}q_{1,N}=0\).  Neither the graph,
its expansion, nor the constants use \(\mathfrak P_{\rm pl}\).
\end{lemma}

\begin{proof}
This is the graph-only result proved by the special-flow construction in
\cref{subsec:invariant-history-details}; specifically, use
\cref{lem:invdet-C0-contraction,prop:invdet-mixed-jets,%
prop:invdet-coefficient-envelope,eq:invdet-local-remainder} with the
retained-tube alternative in that proposition and its local-block proof.
Those arguments assume only the prepared data
\cref{eq:invdet-prepared-data}, which are supplied entirely by
\(\mathfrak P_{\rm gr}\).  The planar trace completion is not defined or
used there.
\end{proof}

\begin{proof}[Proof of \cref{lem:preparation-existence}]
Choose a fixed scalar plateau \(\eta\in C^\infty(\R;[0,1])\), equal to one
on \([-1,1]\) and zero off \((-2,2)\), and a fixed smooth saturation
\(s_{h_*}:\R\to[-2h_*,2h_*]\) equal to the identity on
\([-h_*,h_*]\).  Apply \(s_{h_*}\) componentwise to every stable slot and
then apply \(P_{\perp,N}\).  The product estimate
\cref{eq:shared-products} makes all derivatives of this operation
dimension independent.

In the \((\chi,d)\)-coordinates, the closure of
\(\mathfrak H^{[2]}_\delta\) is compact.  The singular flow is taken only for
\(0\le t\le2\Theta_*\); its variational equations on
\(|\chi|\le S_\delta+1\), \(|d|\le\delta^\kappa\) have derivatives bounded by
\begin{equation}
 C\langle R_\delta\rangle^M e^{cR_\delta}.             \label{eq:preparation-flow-bound}
\end{equation}
In particular,
\(\delta^\kappa e^{cR_\delta}=o(1)\), so fixed
\(B_*>2\Theta_*+B+3\) and \(d_*>b_0\) can be chosen to contain the entire
depth-two agreement hull for all sufficiently small targets.  Use
\cref{eq:fixed-singular-preparation} for the singular field.  Direct
differentiation gives \cref{eq:fixed-singular-lipschitz}; the successive
flow-variational equations give the grade-eleven tame-flow estimate, as
written out in \cref{def:invdet-tame-preparation} and the paragraph following
\cref{eq:invdet-q0-lipschitz}.

For every current and delayed planar slot, rescale the same plateau \(\eta\)
to a rectangle strictly between the preceding agreement hull and the support
of \(\Xi_\delta\).  Apply these plateaus only to the perturbation
nonlinearities in the decomposition
\(q_{0,\delta}+\rho F_{N,\delta}\),
\(A_Nh+\rho G_{N,\delta}\), after the stable saturation already described.
This gives globally bounded \(C^{12}\) maps, agrees with the exact polynomial
chart on every required slot, and uses only the target \(\delta\), not
\((\nu,\zeta)\).  The chain rule, \cref{eq:shared-products}, and
\cref{eq:preparation-flow-bound} give the declared derivative majorant
without a factor depending on \(N\).

It remains to construct the planar completion without rescaling the
shrinking normal width.  Work first in the \((\chi,d)\)-coordinates and put
\[
 K_\delta=[-S_\delta-1,S_\delta+1]
              \times[-\delta^\kappa,\delta^\kappa].
\]
Let \(\widehat Q\) be the global reduced graph field in these coordinates,
let \(\widehat q_0=(1-2\alpha d,\chi d)^T\), and set
\(\widehat f=(\widehat Q-\widehat q_0)|_{K_\delta}\).
By \cref{lem:graph-only-retained-expansion}, for one common \(c,M\), each
mixed jet in \cref{eq:planar-required-jets} satisfies
\begin{equation}
 |D_{\chi,d}^a\partial_\nu^i\partial_\zeta^j\widehat f(\chi,d)|
 \le C\delta^{\omega_{ij}}
       e^{c|\chi|}\langle\chi\rangle^M
 \quad\text{on }K_\delta.                            \label{eq:thin-rectangle-data-bound}
\end{equation}

We use the classical Whitney cube construction
\citep{whitney1934analytic} in a localized weighted form.  Take the smooth
weight
\[
 w(\chi)=e^{c\sqrt{1+\chi^2}}(1+\chi^2)^{M/2};
\]
its first three derivatives, and those of its reciprocal, are bounded by a
constant times \(w\) and \(w^{-1}\), respectively.  For the Whitney cubes
\(\mathcal W_\delta\) of \(\R^2\setminus K_\delta\), choose a nearest point
\(x_C\in K_\delta\) and a subordinate partition
\(\{\phi_C\}_{C\in\mathcal W_\delta}\).  If \(g\in C^3(K_\delta)\), define
\begin{equation}
 W_\delta g=g\ \text{ on }K_\delta,\qquad
 W_\delta g(x)=\sum_{C\in\mathcal W_\delta}
   \phi_C(x)P^3_{x_C}g(x)\ \text{ off }K_\delta,     \label{eq:whitney-thin-extension}
\end{equation}
where \(P^3_{x_C}g\) is the total third-order Taylor polynomial.  Convexity
of \(K_\delta\), Taylor's formula, and the standard Whitney-cube
cancellation give, on every fixed-width neighborhood,
\begin{equation}
 \|W_\delta g\|_{C^3}\le C\|g\|_{C^3(K_\delta)},     \label{eq:whitney-thin-bound}
\end{equation}
with a constant depending only on the dimension, derivative order, and the
chosen fixed collar width, not on the side lengths of \(K_\delta\).  This
construction uses only total state order at most three.

After decreasing \(\delta_{\rm prep}\), if necessary, so that
\(\delta^\kappa<b_0/4\), the separation imposed on \(B\) and \(b_0\)
permits a target-frozen cutoff
\(\varrho_\delta\in C_c^\infty(\R^2;[0,1])\) which equals one on
\(K_\delta\), is supported in a fixed-width collar contained in
\(\{|\chi|<R_\delta-b_0,\ |d|<b_0/2\}\), vanishes on the two endpoint
balls, and satisfies
\(\sup_\delta\|\varrho_\delta\|_{C^3}<\infty\).
Define, componentwise,
\begin{equation}
 E_\delta\widehat f
 =w\,\varrho_\delta
     W_\delta(w^{-1}\widehat f),\qquad
 \widetilde{\widehat Q}
 =q_{0,\delta}^{\chi d}+E_\delta\widehat f.          \label{eq:explicit-planar-completion}
\end{equation}
The operator \(E_\delta\) is linear and target frozen, so it commutes with
\((\nu,\zeta)\)-derivatives.  The bounds in
\cref{eq:thin-rectangle-data-bound,eq:whitney-thin-bound}, the fixed-width
cutoff, and the weight estimates preserve the factor
\(\delta^{\omega_{ij}}\) for every jet in
\cref{eq:planar-required-jets}.

On \(K_\delta\), one has
\(q_{0,\delta}^{\chi d}=\widehat q_0\),
\(\varrho_\delta=1\), and \(E_\delta\widehat f=\widehat f\), so the
completion agrees exactly with the graph field.  On the endpoint balls,
\(E_\delta\widehat f=0\) and
\(q_{0,\delta}^{\chi d}=\widehat q_0\).  On
\(\mathcal C_\delta\), the same equality of the singular fields shows that
\(\widetilde{\widehat Q}-\widehat q_0=E_\delta\widehat f\), which proves
\cref{eq:planar-preparation-smallness}.  Finally,
\(q_{0,\delta}^{\chi d}\) is globally bounded and Lipschitz, while
\(E_\delta\widehat f\) is compactly supported and \(C^3\).
Pullback by the fixed polynomial coordinate map \(\Gamma\) only enlarges
the allowed polynomial weight.  The resulting planar field is globally
bounded, globally Lipschitz, and complete.  All cutoffs, saturation maps,
and Whitney rules are fixed simultaneously for every \(N\), completing the
construction.
\end{proof}

\subsection{Localized obstruction to an intrinsic finite-section root}
\label{subsec:supplementary-root-noncanonicity}

\begin{proof}[Proof of \cref{prop:finite-section-noncanonicity}]
Fix \(N_0\), a sufficiently small target \(\delta_*\), and one planar
completion supplied by \cref{lem:preparation-existence}.  On the attracting
side choose an open strip strictly between the endpoint ball and the retained
physical tube.  The baseline selected trace is embedded there for small
\(\delta_*\), because its \(\chi\)-component is a small perturbation of the
strictly monotone singular parametrization \(\chi=s\).  After narrowing the
strip, the same \(C^1\) control and the separation of the two phase-side
intervals give it positive distance from the repelling trace.

Linearize the attracting trace equation, including its homogeneous endpoint
and phase rows, at the baseline root.  We first show that forcing supported in
the chosen strip has a nonzero central normal response.  Let
\(\Phi_{\delta_*}(s,t)\) be the fundamental matrix of the homogeneous planar
variational equation after the strip.  Since
\(\Phi_{\delta_*}(0,t)\) is invertible, there is a nonzero vector \(y\) at the
right edge \(t\) of the strip for which the first component of
\(\Phi_{\delta_*}(0,t)y\) is zero.  Its second component cannot also vanish.
Choose a smooth variation \(\xi\) which is zero before the strip, joins zero
to \(y\) inside it, and equals
\(\Phi_{\delta_*}(s,t)y\) afterwards.  Writing \(L_{\rm tr}\) for the
linearized trace operator, its residual \(f=L_{\rm tr}\xi\) is supported
inside the artificial strip, its endpoint row and central phase row vanish,
but its central normal component is nonzero.  This is also an immediate
localized consequence of the trace--phase--endpoint isomorphism in
\cref{lem:shared-graph-trace}.

Choose the interpolation so that its residual \(f\) is flat at both
boundaries of the strip.  A tubular extension along the embedded trace
segment then realizes a positive scalar multiple of \(f\) as the restriction
of a compactly supported \(C^\infty\) planar vector-field bump
\(\mathfrak b_*\).  Its support is disjoint from the repelling trace, the
physical agreement tube, and both endpoint balls, and all of its jets vanish
at the support boundary.

It remains to turn this bump at one target into one admissible rule, rather
than choose an unrelated bump for each \(\delta\).  The preceding support has
positive distance from all three forbidden sets.  Hence there is a compact
target interval \(V_*\Subset(0,\delta_{\rm prep})\), containing
\(\delta_*\), on which the same separation persists.  Choose
\(\vartheta_*\in C_c^\infty(V_*;[0,1])\) with
\(\vartheta_*(\delta_*)=1\), and, in the common \((\chi,d)\)-coordinates,
define for every \(N\) and every target \(\delta\)
\begin{equation}
 \widetilde Q_{N,\lambda}
 =\widetilde Q_{N,0}
   +\lambda\vartheta_*(\delta)\mathfrak b_* .         \label{eq:completion-bump-family}
\end{equation}
This is one target-frozen rule family: the added field is independent of
\((\nu,\zeta)\), is the same for all \(N\), and vanishes outside \(V_*\).
Because \(V_*\) is compact and bounded away from zero, a positive rescaling
of \(\mathfrak b_*\) makes all planar mixed-jet preparation bounds in
\cref{eq:planar-required-jets} and
\cref{eq:planar-preparation-smallness} hold with one set of constants.
The rescaling does not destroy the nonzero central response.  Compact support
preserves global boundedness and Lipschitz continuity.  Thus there is a common
\(\lambda_0>0\) for which \cref{eq:completion-bump-family} is admissible for
every \(|\lambda|<\lambda_0\), every target, and every network size.

Differentiate the trace fixed point with respect to \(\lambda\).  The
constructed residual gives a nonzero central normal component, while the
repelling trace is unchanged.  On the phase section,
\(\partial_Y\mathscr H_\alpha=\alpha e/2+o(1)\ne0\); hence
\[
 \left.\partial_\lambda
   \mathscr D^\lambda_{N_0}
      (\delta_*,\nu^{0}_{c,N_0}(\delta_*,0),0)\right|_{\lambda=0}
 \ne0.
\]
The parameterized trace contraction and the finite-dimensional implicit
function theorem used for its endpoint row are jointly \(C^1\) in
\((\nu,\lambda)\) on the common \(\lambda\)-interval.  The root slope at
\(\lambda=0\) is
\(\partial_\nu\mathscr D^0_{N_0}
 =\delta_*\sqrt{2\pi}+O(\delta_*^2)\ne0\).  Applying the implicit-function
theorem with \(\lambda\) as an additional parameter therefore yields
\[
 \left.\partial_\lambda\mu^\lambda_{c,N_0}\right|_{\lambda=0}
 =-\delta_*^2
 \frac{\left.\partial_\lambda\mathscr D^\lambda_{N_0}\right|_{\lambda=0}}
      {\partial_\nu\mathscr D^0_{N_0}}
 \ne0,
\]
which is \cref{eq:completion-root-sensitivity}.
\end{proof}
